\documentclass[12pt]{article}

\usepackage{setspace}
\onehalfspace





\usepackage{amsmath}












































\usepackage{moodle}


\begin{document}

\begin{quiz}{gyakorló}

\begin{multi}[feedback={Egy konstans v.v-nak 0 a szórása.}]{dvv12}
Egy $\xi$ valségi változó szórása mindig pozitív.
\item igaz
\item* hamis
\end{multi}
\begin{multi}[feedback={Nókoment.}]{esem5}
Ha $\mathrm{P}(A)=1$ akkor $A$ bekövetkezhet.
\item* igaz
\item hamis
\end{multi}
\begin{multi}[feedback={Az adott feltételek mellet a várható érték egy súlyozott közép.}]{dvv5}
Egy $\xi$ valségi változó az $x_1,{...\:}, x_n$ értékeket veheti fel. Ekkor:
$$
\min_k(x_k)\le \mathrm{E}(\xi) \le \max_k(x_k)
$$
\item* igaz
\item hamis
\end{multi}
\begin{multi}[feedback={Def.}]{impfvv11}
Ha $\xi\sim {\cal U}(a,b)$ akkor a sűrűségfüggvénye 
$$
f(x)=
\begin{cases}
\frac{1}{b-a} & \ \ a<x<b \\
0 & \text{máskor} 
\end{cases}
$$
\item* igaz
\item hamis
\end{multi}
\begin{multi}[feedback={Tétel.}]{flenvv3}
Ha $\xi$ és $\eta$ valségi változók függetlenek, akkor
$$
\mathrm{E}(\xi \eta)=\mathrm{E}(\xi)\mathrm{E}(\eta)
$$
feltéve, hogy létezik a szórásuk.
\item* igaz
\item hamis
\end{multi}
\begin{multi}[feedback={Igen, éppen ez alapján tudunk dönteni. Ahhoz képest hogy sztenderd normális átlagos, tipikus értéket, vagy kiugró, extrém értéket látunk.}]{proba9}
Ha egy $16$ elemű minta normális sokaságból származik, $\sigma=3$ és hipotéziseink:
$$
H_0: \mu=10
$$
$$
H_1: \mu\neq 10
$$
akkor az 
$$
\frac{\overline{X}-10}{3}4
$$
valségi változó sztenderd normális eloszlású ha $H_0$ igaz.
\item* igaz
\item hamis
\end{multi}
\begin{multi}[feedback={Def.}]{impdvv5}
Azt mondjuk, hogy $\xi\sim \mathrm{Binom}(n,p)$, ha 
$$
\mathrm{P}(\xi=k)=\binom{n}{k}p^k(1-p)^{n-k}\ \ \ k=0{...\:} n
$$
valamilyen $n>0$ és $p\in[0,1]$ esetén.
\item* igaz
\item hamis
\end{multi}
\begin{multi}[feedback={Definíció}]{minta2}
Egy $X_1,{...\:},X_n\sim X$ minta esetén a mintaátlag $\overline{X}=\frac{X_1+{...\:}+X_n}{n}$
\item* igaz
\item hamis
\end{multi}
\begin{multi}[feedback={Például: legyen a kísérlet random húzás a $[0,1]$-ből és $A=\{x\in [0,1]:x\neq \frac{1}{2}\}$}]{esem6}
Ha $\mathrm{P}(A)=1$ akkor $A$ mindig bekövetkezik.
\item igaz
\item* hamis
\end{multi}
\begin{multi}[feedback={$A=B$?}]{val3}
A valószínűség monoton, azaz $\mathrm{P}(A)<\mathrm{P}(B)$, ha $A\subseteq B$.
\item igaz
\item* hamis
\end{multi}
\begin{multi}[feedback={Egyenletes lineáris transzformáltja is egyenletes (a konstans elfajult esettől eltekintve)}]{impfvv5}
Ha $\xi\sim {\cal U}(0,1)$ akkor minden $a<b$-re $(b-a)\xi+a\sim {\cal U}(a,b)$.
\item* igaz
\item hamis
\end{multi}
\begin{multi}[feedback={A páronkénti gyengébb tulajdonság.}]{felt8}
A páronkénti függetlenségből következik a teljes.
\item igaz
\item* hamis
\end{multi}
\begin{multi}[feedback={Definíció}]{minta3}
Egy $X_1,{...\:},X_n\sim X$ minta esetén a mintaátlag $\overline{X}=\frac{X_1+{...\:}+X_n}{n-1}$
\item igaz
\item* hamis
\end{multi}
\begin{multi}[feedback={A fenti formula az $t$-próba ún. próbastatisztikája.}]{proba7}
Az $u$-próbánál a 
$$
u=\frac{\overline{X}-\mu_0}{s}\sqrt{n}
$$
mennyiség segítségével döntünk a hipotéziseinkről.
\item igaz
\item* hamis
\end{multi}
\begin{multi}[feedback={wiki}]{impfvv1}
Egy $\xi\sim {\cal U}(a,b)$ várható értéke: $\frac{a+b}{2}$
\item* igaz
\item hamis
\end{multi}
\begin{multi}[feedback={$b\xi+a \sim {\cal U}(a,a+b)$}]{impfvv6}
Ha $\xi\sim {\cal U}(0,1)$ akkor  minden $a<b$-re $b\xi+a\sim {\cal U}(a,b)$.
\item igaz
\item* hamis
\end{multi}
\begin{multi}[feedback={de-Morgan azonosság}]{esem3}
Az $\overline{A\cdot B}$ maga után vonja az $\overline{A} + \overline{B}$ eseményt.
\item* igaz
\item hamis
\end{multi}
\begin{multi}[feedback={Def.}]{dvv9}
Egy $\xi$ valségi változó az $x_1,{...\:}, x_n$ értékeket veheti fel és $p_k=\mathrm{P}(\xi=x_k)$. 
Ekkor a második momentuma:
$$
\mathrm{E}(\xi^2)=\sum_{k=1}^n x_k p_k^2
$$
\item igaz
\item* hamis
\end{multi}
\begin{multi}[feedback={Definíció.}]{esem8}
Az $A$ és $B$ események pont akkor függetlenek, ha $\mathrm{P}(A+B)=\mathrm{P}(A)+\mathrm{P}(B)$.
\item igaz
\item* hamis
\end{multi}
\begin{multi}[feedback={A fenti formula az $u$-próba ún. próbastatisztikája.}]{proba5}
Az $u$-próbánál a 
$$
u=\frac{\overline{X}-\mu_0}{\sigma}\sqrt{n}
$$
mennyiség segítségével döntünk a hipotéziseinkről.
\item* igaz
\item hamis
\end{multi}
\begin{multi}[feedback={Additívitás.}]{val6}
Bármely esemény valsége kiszámolható a $\frac{\text{kedvező}}{\text{összes}}$ képlettel, ha az elemi események 
egyforma valségűek.
\item* igaz
\item hamis
\end{multi}
\begin{multi}[feedback={Tulajdonság,tétel.}]{flenvv9}
Ha a $\xi$ és $\eta$ valségi változóknak létezik a korrelációja, akkor $-1\le \mathrm{corr}(\xi,\eta)\le 1$.
\item* igaz
\item hamis
\end{multi}
\begin{multi}[feedback={Csak ha az elemi események egyforma valségűek.}]{val5}
Bármely esemény valsége kiszámolható a $\frac{\text{kedvező}}{\text{összes}}$ képlettel.
\item igaz
\item* hamis
\end{multi}
\begin{multi}[feedback={Ebből nem jön ki a pl. a monotonitás vagy a balról folytonosság.}]{eofv2}
Azt mondjuk, hogy egy $F:\mathbb{R}\to \mathbb{R}$ eloszlásfüggvény, ha $F\ge 0$ és 
$$
\int_{-\infty}^{+\infty}F(x)\mathrm{d } x=1
$$
\item igaz
\item* hamis
\end{multi}
\begin{multi}[feedback={$\{a\le\xi<b\}=\{\xi<b\}\setminus \{\xi<a\}$}]{eofv7}
Bármely $\xi$-re és $a<b$-re $\mathrm{P}(a\le \xi <b)=F_{\xi}(b)-F_{\xi}(a)$.
\item* igaz
\item hamis
\end{multi}
\begin{multi}[feedback={Def.}]{flenvv1}
Azt mondjuk, hogy $\xi$ és $\eta$ diszkrét valségi változók függetlenek, ha
$$
\mathrm{P}(\xi=x,\eta=y)=\mathrm{P}(\xi=x)\mathrm{P}(\eta=y)
$$
bármely $x,y\in \mathbb{R}$-re teljesül.
\item* igaz
\item hamis
\end{multi}
\begin{multi}[feedback={Ezek az $u$-próba használatának előzményei.}]{proba1}
Ha a minta normális sokaságból származik, ismert a szórás és a sokasági várható értékre (átlagra) vonatkozó kérdésről 
akarunk dönteni akkor $u$-próbát használhatunk.
\item* igaz
\item hamis
\end{multi}
\begin{multi}[feedback={$\xi=\xi_1+{...\:}+\xi_n$, ahol $\xi_k$ Bernoulli és $\mathrm{D}^2(\xi_k)=p(1-p)$+$\mathrm{D}^2$-additívitása(függetlenség esetén)}]{impdvv4}
Egy $\xi\sim \mathrm{Binom}(n,p)$ szórásnégyzete: $\frac{n+1}{2}$
\item igaz
\item* hamis
\end{multi}
\begin{multi}[feedback={Az $s^2$ torzítatlan becslése a sokasági (elméleti) szórásnégyzetnek}]{minta7}
Az $s^2$ korrigált empirikus szórásnégyzetre $\mathrm{E}(s^2)=\mathrm{D}^2(X)$.
\item* igaz
\item hamis
\end{multi}
\begin{multi}[feedback={Diszkrét v.v.-ra ez nem feltétlen igaz, ha $a$-t pozitív valséggel felveszi.}]{eofv8}
Bármely $\xi$-re és $a<b$-re $\mathrm{P}(a< \xi <b)=F_{\xi}(b)-F_{\xi}(a)$.
\item igaz
\item* hamis
\end{multi}
\begin{multi}[feedback={Def.}]{impfvv8}
Ha $\xi\sim {\cal N}(0,1)$ akkor a sűrűségfüggvénye $f(x)=\frac{e^{-\frac{x^2}{2}}}{\sqrt{2\pi}}$.
\item* igaz
\item hamis
\end{multi}
\begin{multi}[feedback={Komplementer események.}]{eofv9}
Bármely $\xi$-re és $a$-ra $F_{\xi}(a)+\mathrm{P}(\xi\ge a)=1$.
\item* igaz
\item hamis
\end{multi}
\begin{multi}[feedback={Def.}]{dvv6}
Egy $\xi$ valségi változó az $x_1,{...\:}, x_n$ értékeket veheti fel. Ekkor a szórása:
$$
\mathrm{D}(\xi)=\sum_{k=1}^n \mathrm{P}(\xi=x_k)(x_k-\mathrm{E}(\xi))^2
$$
\item igaz
\item* hamis
\end{multi}
\begin{multi}[feedback={Def.}]{felt4}
$A_1,{...\:},A_n$ teljes eseményrendszer, ha $\sum A_k=\Omega$ és a tagok páronként függetlenek.
\item igaz
\item* hamis
\end{multi}
\begin{multi}[feedback={Def.}]{flenvv4}
Ha $\xi$ és $\eta$ diszkrét valségi változóknak létezik a szórása és nemnulla, akkor a 
korrelációjuk
$$
\mathrm{corr}(\xi,\eta)=\frac{\mathrm{cov}(\xi,\eta)}{\mathrm{D}(\xi)\mathrm{D}(\eta)}
$$
\item* igaz
\item hamis
\end{multi}
\begin{multi}[feedback={A feltételek mellett a fenti mennyiség sztenderd normális, ha $H_0$ igaz.}]{proba10}
Ha egy $16$ elemű minta normális sokaságból származik, $\sigma=3$ és hipotéziseink:
$$
H_0: \mu=10
$$
$$
H_1: \mu\neq 10
$$
akkor az 
$$
\frac{\overline{X}-10}{3}4
$$
valségi változó $15$-szabadsági fokú Student eloszlású ha $H_0$ igaz.
\item igaz
\item* hamis
\end{multi}
\begin{multi}[feedback={Def.}]{dvv7}
Egy $\xi$ valségi változó az $x_1,{...\:}, x_n$ értékeket veheti fel. Ekkor a szórásnégyzete:
$$
\mathrm{D}^2(\xi)=\sum_{k=1}^n \mathrm{P}(\xi=x_k)(x_k-\mathrm{E}(\xi))^2
$$
\item* igaz
\item hamis
\end{multi}
\begin{multi}[feedback={Def.}]{impdvv8}
Azt mondjuk, hogy $\xi\sim \mathrm{Poisson}(\lambda)$, ha 
$$
\mathrm{P}(\xi=k)=e^{\lambda}\frac{\lambda^k}{k}\ \ \ (k=0,1,2{...\:} )
$$ 
valamilyen $\lambda>0$ esetén.
\item igaz
\item* hamis
\end{multi}
\begin{multi}[feedback={Def.+$P(AB)=P(A)$ a feltételekből.}]{felt1}
$A$-nak a $B$-re ($\mathrm{P}(B)>0$) vonatkozó feltételes valsége $\mathrm{P}(A|B)=\frac{\mathrm{P}(A)}{\mathrm{P}(B)}$, ha $A$ maga után vonja $B$-t.
\item* igaz
\item hamis
\end{multi}
\begin{multi}[feedback={$f\ge 0$ fontos (monoton nemcsökkenő $F$-nek a deriváltja)}]{sfv3}
Azt mondjuk, hogy egy $f:\mathbb{R}\to \mathbb{R}$ sűrűségfüggvény 
$$
\int_{-\infty}^{+\infty}f(x)\mathrm{d } x\ge 0
$$
\item igaz
\item* hamis
\end{multi}
\begin{multi}[feedback={Def.}]{sfv7}
Ha $\xi$-nek $f$ a sűrűségfüggvénye, akkor bármely $a<b$-re:
$$
\mathrm{P}(a<\xi<b)=f(b)-f(a)
$$
\item igaz
\item* hamis
\end{multi}
\begin{multi}[feedback={Def.}]{sfv6}
Ha $\xi$ sűrűségfüggvénye $f$, akkor bármely $a<b$-re:
$$
\mathrm{P}(a<\xi<b)=\int_{a}^{b}f(x)\mathrm{d } x
$$
\item* igaz
\item hamis
\end{multi}
\begin{multi}[feedback={Def.}]{sfv2}
Egy $f:\mathbb{R}\to [0,1]$ pontosan akkor sűrűségfüggvény, ha monoton nemcsökkenő, 
balról folytonos, $\lim_{x\to+\infty}f(x)=1$ és $\lim_{x\to-\infty}f(x)=0$
\item igaz
\item* hamis
\end{multi}
\begin{multi}[feedback={Tétel. Ez tetszőleges v.v-kre is igaz.}]{dvv15}
$\xi$ és $\eta$ $\xi$ és $\eta$ véges értékkészletű valségi változókra:
$$
\mathrm{D}^2(\xi+\eta)=\mathrm{D}^2(\xi)+\mathrm{D}^2(\eta)+2\mathrm{cov}(\xi,\eta)
$$
\item* igaz
\item hamis
\end{multi}
\end{quiz}

\end{document}

